\uuid{cLCp}
\exo7id{1383}
\titre{exo7 1383}
\auteur{ortiz}
\organisation{exo7}
\datecreate{1999-04-01}
\isIndication{false}
\isCorrection{false}
\chapitre{Groupe, anneau, corps}
\sousChapitre{Algèbre, corps}
\cpge{mpsi-mp2i.structures-algebriques.c, mpsi-mp2i.espaces-vectoriels.a, pcsi.espaces-vectoriels.a, ptsi.geometrie-espace.c, mpsi-mp2i.applications-lineaires.b, pcsi.applications-lineaires.b, ptsi.applications-lineaires.b, mp-mpi.reduction-endomorphismes.g, pc.complements-espaces-vectoriels.d, psi.complements-espaces-vectoriels.d, bcpst2.applications-lineaires-matrices.a}
\module{Algèbre}
\niveau{L2}
\difficulte{}

\contenu{
\question{
Soient $E$ un $\Rr$-ev et $u\in \mathcal{L} (E)$
 tel que $u^2=u$. On d\'efinit
$${\Rr[u]=\left\{a\text{Id}_E+bu:a,b\in \Rr\right\}.}$$
 Montrer que,
muni des lois usuelles sur $\mathcal{L}(E)$, c'est
une $\Rr$-alg\`ebre. L'anneau $\Rr[u]$ est-il un
corps ?
}
}
